\section{Threat Model and Field-Level Risk}\label{sec:problem}\label{sec:mk}
%===============================================================================

\subsection{Parties, field roles and threat model}
The \textbf{auditor} is the data owner and evaluates confidentiality before release. It holds the complete history of all operational
fields, including those it regards as sensitive, and grades fields offline before release.
The \textbf{adversary} obtains published fields and is assumed to hold historical auxiliary
labels of the sensitive target, from which it trains a supervised inference model; risk is
therefore measured as out-of-sample \emph{supervised} inferability. The adversary may combine
the field under review with other fields it already holds.

Field roles are fixed by disclosure rules, not by data analysis. Let $G$ be the candidate
fields that the rules allow to be published, $B\subseteq G$ a \emph{public baseline} that is
already released ($B=\emptyset$ unless stated otherwise), $H=G\setminus B$ the fields under
review, and $Y_y$ a sensitive target. Section~\ref{sec:data} instantiates these roles for
four data sets from the categories and release timings of the Chinese electricity-market
disclosure rules~\cite{nea2024disclosure}. An ex-ante estimate of the \emph{same physical
quantity} as a target (e.g.\ a load forecast when the target is metered load) is removed from
$G$, because it is a copy of the target rather than an inference channel.

\subsection{Set-level inferability}
\begin{definition}[Population inferability]
\[
\Vy^{*}(S) = 1 - \inf_{f}\,\mathbb{E}\big[(Y_y - f(X_S))^{2}\big] \big/ \mathrm{Var}(Y_y),
\]
the infimum taken over measurable $f$. It satisfies $\Vy^{*}(\emptyset)=0$, $0 \le \Vy^{*} \le 1$
and is monotone in $S$.
\end{definition}

\begin{definition}[Empirical inferability]\label{def:vret}
The data are split into training and test rows, with a validation part inside training. For
an attacker family $\mathcal{H}$, every $h\in\mathcal{H}$ is trained on $X_S$ only; the
attacker and its training epoch are selected on validation, and the test $\Rtwo$ is reported:
$\Vy^{\mathrm{ret}}(S) = \Rtwo_{\mathrm{test}}(\hat h_S)$ with
$\hat h_S = \argmax_{h \in \mathcal{H},\,\mathrm{epoch}} \Rtwo_{\mathrm{val}}(h;S)$.
\end{definition}

An adversary holding a set can also use any of its subsets. We therefore evaluate the subset-closed inferability
\begin{equation}\label{eq:closure}
\Vy(S)=\max_{T\subseteq S}\big[\Vy^{\mathrm{ret}}(T)\big]_+,\qquad \Vy(\emptyset)=0.
\end{equation}
This retraining-based reference defines the empirical risk tables used throughout the paper.

\subsection{Contextual gain}
The quantity from which everything else follows is the gain that releasing field $i$ brings over a
given background $T$ of fields the adversary already holds,
\begin{equation}\label{eq:delta}
\Delta_i(A)=\Vy(A \cup \{i\}) - \Vy(A),\qquad A=B\cup T .
\end{equation}
$\Delta_i(B)$ is the single-field risk. The same field can be nearly useless over one background
and decisive over another: a price series says little about a unit's output alone, but together
with the unit's forecast availability it reveals the residual quantity the unit covers. A
field-level risk quantity must therefore summarize a whole \emph{distribution} of contextual gains
$\{\Delta_i(B\cup T)\}_T$, and the question is which summary serves security.

\subsection{Budgeted marginal risk and background amplification}
\begin{definition}[Budgeted marginal risk, background amplification]\label{def:mk}
For target $y$, public baseline $B$, fields under review $H=G\setminus B$ and budget $K$,
\begin{align}
\Mik &= \max_{T \subseteq H \setminus \{i\},\, |T| \le K}\ \Delta_i(B\cup T), \label{eq:mk}\\
\Gamma^{(K)}_{i\to y} &= \Mik - M^{(0)}_{i\to y} \;\ge 0 . \label{eq:gamma}
\end{align}
Any maximizer $T^{*}$ is a \emph{witness background}.
\end{definition}
$\Mik$ is the largest increase of inferability that releasing $i$ can cause when the adversary
holds at most $K$ other fields; $\Gamma^{(K)}_{i\to y}$ is the part of it that exists only because of
the background, i.e.\ the amplification of single-field risk by combination. For each (field,
target) pair the audit reports the profile $\big(M^{(0)}_{i\to y},\ \Gamma^{(K)}_{i\to y},\ \Mik\big)$
together with the top-$r$ backgrounds ranked by $\Delta_i$; in our data the top backgrounds are
often nearly tied and share a core partner field, so the partner is more informative than a
single argmax. At $K=1$ the amplification has a closed form.

\begin{proposition}[Amplification at $K=1$]\label{prop:k1}
With $I^B_{ij}=\Vy(B\cup\{i,j\})-\Vy(B\cup\{i\})-\Vy(B\cup\{j\})+\Vy(B)$,
$\Gamma^{(1)}_{i\to y}=\max\big(0,\max_{j} I^{B}_{ij}\big)$.
\end{proposition}
\begin{proof}
Over the background $\{j\}$, $\Delta_i(B\cup\{j\})=M^{(0)}_{i\to y}+I^B_{ij}$; over $\emptyset$ it is
$M^{(0)}_{i\to y}$. Take the maximum.
\end{proof}
A field is amplified by one background field exactly when it is strictly super-additive with it,
and the amplification equals the strongest positive interaction.

\subsection{Relation to feature importance}
An average of contextual gains, as in Shapley values or SAGE~\cite{covert2020sage}, describes typical contribution. The maximum asks whether an admissible background produces a large gain, which is the relevant question when reviewing a release. In our data the median ratio of the average gain to this maximum is 0.32 (Section~\ref{sec:res-amp}). The MCI~\cite{catav2021mci} formalizes this maximum contribution; the public baseline and background budget specialize it to the audit setting. \ref{app:requirements} compares the requirements met by alternative field scores.

\begin{proposition}[Relation to MCI]\label{prop:mci}
Let $v_B(T) := \Vy(B \cup T) - \Vy(B)$ for $T \subseteq H$. Then $v_B$ is monotone,
non-negative and $v_B(\emptyset)=0$; $\Mik = \max_{|T|\le K}[v_B(T\cup\{i\}) - v_B(T)]$ is the
budget-constrained MCI of $v_B$; for $K \ge |H|-1$ and $B=\emptyset$ it is the MCI of $\Vy$.
\end{proposition}
The axiomatic properties of MCI for monotone evaluation functions~\cite{catav2021mci}---dummy,
symmetry, the self-contribution lower bound $\Mik\ge M^{(0)}_{i\to y}$ and invariance to
duplicated features---therefore transfer and are cited rather than claimed; $\Mik$ is
non-decreasing in $K$ and reaches the MCI at $K=|H|-1$.

\subsection{Security meaning}
\begin{theorem}[Threshold-free safety margin]\label{thm:margin}
Fix $i$, $y$, $B$, $K$. A number $m\ge 0$ satisfies ``for every threshold $\tau$ and every
background $|T|\le K$: $\Vy(B\cup T)\le\tau-m \Rightarrow \Vy(B\cup T\cup\{i\})\le\tau$''
if and only if $m\ge\Mik$. For a fixed $\tau$, the smallest $m$ with this property can be
strictly smaller than $\Mik$.
\end{theorem}
\begin{proof}
If $m\ge\Mik$ then $\Vy(B\cup T\cup\{i\})=\Vy(B\cup T)+\Delta_i(B\cup T)\le(\tau-m)+m$.
Conversely, for any $T$ choose $\tau=\Vy(B\cup T)+m$; the property gives $\Delta_i(B\cup T)\le
m$. Strictness: with $B=\emptyset$, $K=1$, $\Vy(\{i\})=0.1$, $\Vy(\{j\})=0.5$,
$\Vy(\{i,j\})=0.9$ one has $M^{(1)}_{i\to y}=0.4$, yet at $\tau=0.95$ the necessary margin is $0$.
\end{proof}
Thus $\Mik$ is the answer to ``by how much can publishing $i$ raise the adversary's accuracy,
whatever it already holds within budget''---a property of the field that does not depend on
the threshold later chosen for grading.

\begin{definition}[Criticality and minimal unsafe sets]
Under baseline $B$, $E\subseteq H$ is a $\tau$-\emph{minimal unsafe set} (MUS) if
$\Vy(B\cup E)>\tau$ and $\Vy(B\cup E')\le\tau$ for every $E'\subsetneq E$. Field $i$ is
$K$-\emph{background} $\tau$-\emph{critical} if some $T$ with $|T|\le K$, $i\notin T$, has
$\Vy(B\cup T)\le\tau<\Vy(B\cup T\cup\{i\})$.
\end{definition}

\begin{theorem}[Criticality $\iff$ small MUS]\label{thm:mus}
If $\Vy$ is monotone and $\Vy(B)\le\tau$, field $i$ is $K$-background $\tau$-critical if and
only if it belongs to a MUS of size at most $K+1$.
\end{theorem}
\begin{proof}
($\Leftarrow$) For a MUS $E\ni i$ with $|E|\le K+1$ take $T=E\setminus\{i\}$.
($\Rightarrow$) If $\Vy(B\cup T)\le\tau<\Vy(B\cup T\cup\{i\})$, take an inclusion-minimal unsafe
$E\subseteq T\cup\{i\}$; it is a MUS of size $\le K+1$. If $i\notin E$ then $E\subseteq T$ and
monotonicity gives $\Vy(B\cup T)\ge\Vy(B\cup E)>\tau$, a contradiction.
\end{proof}
Theorem~\ref{thm:mus} makes the field-level flag a faithful projection of the set-level
danger structure: flagging the $\tau$-critical fields misses no field that takes part in a
dangerous combination of at most $K+1$ fields, and a threshold-crossing background identifies the combination.
The same list of minimal unsafe sets also supports withholding decisions
(Section~\ref{sec:application}).



\subsection{Background budget}
The budget $K$ specifies how many additional fields the adversary may combine with the field under review, beyond the public baseline. We use $K=2$ for the main audit and evaluate $K=3$ on every data set. Small budgets capture interpretable combinations and control the enumeration cost. Section~\ref{sec:res-robust} reports both changes in marginal risk and changes in threshold crossings as the budget increases.
